\section*{v1.9.0 Readout Stabilization and Proof Memory Core}
\addcontentsline{toc}{section}{v1.9.0 Readout Stabilization and Proof Memory Core}
\begin{quote}
\textbf{Normalization note.} This block is retained as a historical precursor of the normalized readout architecture. Its proof-bearing content is re-exported later through the grouped-layer normalization, especially the readout/collapse phases and the semantic--tensor transport discipline. It should therefore be read as historical proof memory rather than as an independent new foundation.
\end{quote}

This release adds a further stabilization layer to the inherited closure, probe,
projection, and appendix architecture. Its purpose is to make the final readout
path more disciplined: not only must admissible content survive transport, it
must also remain recoverable after compression. In this sense v1.9.0 pushes
the document from a transport-stable monolith toward a proof-memory stable one.
A statement is not treated as mature merely because it can be compressed into a
lemma, table row, or projection line; it must also remain reconstructible as the
same closure-supported content when the compression is unfolded again.

\subsection*{1. Readout stabilization principle}
The main strengthening of v1.9.0 is that every admissible readout must be
stable under re-expansion. A readout is the visible endpoint of a structural
chain, but it remains acceptable only if its origin can be reopened without
changing the admissibility class. Thus the document now treats readout as a
\emph{controlled fixed point} of compression and reconstruction:
\[
\begin{aligned}
&\text{local structure} \Longrightarrow \text{HC} \Longrightarrow \text{closure}\\
&\Longrightarrow \text{probe filtration} \Longrightarrow \text{table/tool compression}\\
&\Longrightarrow \text{readout} \Longrightarrow \text{re-expansion to the same admissibility class}.
\end{aligned}
\]

\paragraph{Stabilization Proposition A (readout fixed-point discipline).}
If a statement is compressed into a readout form by admissibility-preserving
transports and can be re-expanded to the same closure-licensed content, then the
readout is structurally stable. If re-expansion changes the admissibility class,
the readout was only heuristic.

\paragraph{Proof sketch.}
The admissibility source is fixed upstream by the ground condition, HC, and
closure. Compression is legitimate only if it does not sever that source. Hence
re-expansion must recover the same structural ancestry rather than a merely
similar narrative.

\subsection*{2. Tool-to-lemma synchronization}
Previous versions upgraded tools and tables into proof interfaces. v1.9.0
adds a stricter synchronization rule: every reusable tool should point to at
least one lemma-level justification, and every repeated lemma pattern should be
compressible into a named tool or interface move once stabilized. The goal is to
prevent the appendix from splitting into disconnected ``tricks'' on one side and
formal arguments on the other.

\paragraph{Stabilization Proposition B (tool/lemma synchronization).}
A reusable tool is mature only when it is synchronized with a lemma pattern that
licenses its use, and a recurring lemma pattern is fully normalized only when it
admits a stable compressed tool form.

\paragraph{Proof sketch.}
Tools without lemma ancestry become opaque shortcuts; lemmas without compressed
interface forms bloat the running line unnecessarily. Synchronization preserves
both transparency and speed while keeping the proof architecture auditable.

\subsection*{3. Probe-to-readout compatibility}
The probe channel now receives a further role. It is not only a gate before
projection but also a compatibility filter for later readouts. Any sector that
passes through probe filtration should emerge with enough structural trace to be
placed into a table, a compressed lemma, or a projection sentence without losing
its closure ancestry.

\paragraph{Stabilization Proposition C (probe-filtered modes carry readout ancestry).}
If a mode survives the probe filtration under the standing orthogonality and HC
rules, then its later readout must retain a backward trace to the probe channel
that selected it. Readout therefore cannot erase probe ancestry without lowering
its status to heuristic interpretation.

\subsection*{4. Appendix as proof memory lattice}
The appendix is sharpened once more: it is now treated as a \emph{proof memory
lattice}. Its entries are not merely stored results but nodes from which one can
move upward to the local closure chain, sideways to equivalent interface forms,
and downward to bounded projection language. This lattice viewpoint organizes the
appendix as a recoverability structure rather than a list of leftovers.

\paragraph{Stabilization Principle (memory-lattice admissibility).}
An appendix entry is fully admissible only if it supports three motions:
\begin{enumerate}[leftmargin=2em,label=(\alph*)]
  \item upward recovery to its closure/lemma ancestry,
  \item lateral transport to equivalent tool or table interfaces, and
  \item downward use in a bounded readout or projection context.
\end{enumerate}
If one of these motions fails, the entry must be marked as provisional,
heuristic, or not yet fully normalized.

\subsection*{5. Operational consequence for later versions}
With v1.9.0 the monolith receives a clearer maturity criterion for future
expansions. New derivations should not merely add content; they should improve
recoverability. In practice this means that later mathematical deepening should
prefer additions that tighten one of the following links: closure to probe,
probe to table, table to lemma, lemma to tool, or readout back to structural
ancestry. This keeps the entire project aligned with the structure-first rule
while making the growing proof apparatus easier to audit and reuse.

\paragraph{Operational rule for subsequent releases.}
Whenever a new argument is introduced, one should ask not only whether the step
is locally correct, but also whether it improves the global recoverability of
compressed proof content. The strongest additions are those that make the same
statement simultaneously easier to compress, easier to re-expand, and easier to
locate inside the appendix memory lattice.



\section*{v2.0.0 Structural Closure Ledger and Expansion Discipline Core}
\addcontentsline{toc}{section}{v2.0.0 Structural Closure Ledger and Expansion Discipline Core}
\begin{quote}
\textbf{Normalization note.} This block is retained as a historical ledger/governance precursor. In the normalized architecture its active role belongs chiefly to the projection/translation layer and to the later formal unfolding phases that carry the actual proof-bearing closure and transport statements.
\end{quote}

This release upgrades the monolith from a proof-memory stable object to a more
explicit \emph{closure ledger}. The point is not merely to preserve ancestry in
principle, but to keep a disciplined account of how every compressed statement,
interface row, or projection sentence remains licensed by the same upstream
closure architecture. In this sense v2.0.0 turns recoverability into an
operational bookkeeping rule: admissible content should be traceable, expandable,
and relocatable without silently changing its structural status.

\subsection*{1. Closure ledger principle}
A closure ledger is the ordered record of the transports that make a statement
usable. It does not store every intermediate calculation in full; rather, it
stores the \emph{licensed path class} by which the content travels from local
structure to admissible readout. The relevant chain is now treated as a ledgered
sequence:
\[
\begin{aligned}
&\text{local triolic structure} \Longrightarrow \text{ground condition} \Longrightarrow \text{HC compensation}\\
&\Longrightarrow \text{closure confinement} \Longrightarrow \text{probe/interface filtration}\\
&\Longrightarrow \text{lemma or tool compression} \Longrightarrow \text{bounded projection/readout}.
\end{aligned}
\]
A later re-expansion is admissible only if it re-enters the same ledger class,
not merely a similar verbal narrative.

\paragraph{Ledger Proposition A (transport accountancy).}
A compressed claim is structurally mature only if one can specify the closure
ledger that licenses its transport. If the path cannot be stated, the claim may
still be suggestive, but it has not yet acquired full structural status.

\paragraph{Proof sketch.}
The earlier versions already required closure ancestry and re-expandability. The
present strengthening adds explicit accountancy: the transport steps must be
stable enough to be named, tracked, and checked. This prevents hidden leaps from
entering the proof line under the cover of familiar notation.

\subsection*{2. Expansion discipline as a release criterion}
The next layer of consolidation concerns how new material may be added. A later
expansion should no longer be judged only by local correctness or plausibility;
it should also improve the ledger quality of the monolith. Thus expansion is now
classified by whether it strengthens one of the major controlled passages:
closure to probe, probe to interface, interface to lemma, lemma to readout, or
readout back to closure ancestry.

\paragraph{Ledger Proposition B (expansion must improve auditability).}
A new derivation is preferable to an older one when it makes the same structural
content easier to audit across compression levels. Additions that only increase
volume without sharpening one of the ledger passages remain secondary.

\paragraph{Operational reading.}
This gives the document a stricter maturity test. Future versions should not add
formal bulk merely because it is available; they should add material that makes
more of the proof architecture checkable under controlled compression and
re-expansion.

\subsection*{3. Probe-ledger compatibility}
The probe channel is now interpreted as a \emph{ledger gate}. It does not simply
filter modes before projection; it also determines which later compressed forms
inherit admissibility in a controllable way. A mode that survives probe
filtration under orthogonality and HC should arrive downstream with a stable
ledger signature that can be recognized in lemma, table, and readout form.

\paragraph{Ledger Proposition C (probe selection induces ledger signature).}
If a mode is selected by the admissible probe channel, then any later compressed
representation of that mode should preserve a recognizable signature of that
selection. If the signature disappears completely, the downstream statement must
be marked as heuristic or only partially normalized.

\subsection*{4. Table and appendix discipline}
The appendix and its tables now receive a slightly sharper role again. They are
not only memory lattices; they function as \emph{ledger surfaces}. Each row,
entry, or named tool should indicate which closure path class it abbreviates and
which re-expansion channel it expects. This gives the appendix a more active
status inside the proof architecture and reduces the danger that tables drift
into mere summaries.

\paragraph{Ledger Proposition D (rows abbreviate path classes).}
A normalized table row is admissible only if it abbreviates a stable path class
rather than an isolated sentence. The row then becomes reversible in two senses:
it can be expanded upward into lemma ancestry and laterally into equivalent tool
or interface forms.

\subsection*{5. Transition to the next mathematical deepening}
With v2.0.0 the project crosses a useful threshold. The document now has a
more explicit rule for how future mathematical elaboration should be inserted:
not as detached calculations, but as ledger-preserving thickening of already
licensed paths. This is especially relevant for the larger companion roadmap,
where many later lemmas and proof tools must remain structurally synchronized.

\paragraph{Working rule for subsequent versions.}
Whenever a new proof block is added, it should be possible to say which closure
ledger it refines, which compressed interfaces it improves, and how its content
returns to the appendix without losing admissibility. Material that satisfies
all three demands strengthens the monolith in the intended structure-first way.


\clearpage
\section*{v2.1.0 Structural Audit and Re-Expansion Core}
\addcontentsline{toc}{section}{v2.1.0 Structural Audit and Re-Expansion Core}
\begin{quote}
\textbf{Normalization note.} This block is retained as a historical audit precursor. Its main value is to preserve re-expansion discipline and proof-memory traceability; the normalized document no longer treats it as a separate generative core, but as supporting governance for the later unfolding phases.
\end{quote}

This release adds a further consolidation layer to the companion monolith. The
new emphasis is not merely on carrying closure information forward, but on
making explicit how later re-expansion remains auditable after several stages of
compression. Earlier versions already introduced closure ladders, interface
filters, ledger surfaces, and proof-memory discipline. Version~v2.1.0 now
adds an \emph{audit core}: every admissible compressed statement should retain a
recoverable expansion route that can be checked against the same closure class
from which it descended.

\subsection*{1. Auditability as structural maturity}
A compressed statement becomes structurally mature only when it is possible to
state not just where it came from, but how it may be unfolded again without
changing its admissible class. This requirement sharpens the earlier ledger
ideas. The ledger tracked which closure path licensed a statement; the audit
core now asks whether that path can be re-opened in a controlled way.

\paragraph{Audit Proposition A (licensed re-expansion).}
A normalized claim is fully admissible only if there exists a bounded
re-expansion route from the compressed form back to a compatible lemma, tool,
probe, or closure ancestry class. A statement that cannot be re-opened under
controlled transport may still be heuristically useful, but it has not yet
reached final structural normalization.

\paragraph{Proof sketch.}
Closure ancestry alone prevents arbitrary insertion of unsupported projection
claims, but it does not yet guarantee that later readers can recover the path.
The present strengthening adds a re-opening condition. The compressed surface
must not merely summarize the proof memory; it must preserve enough structure
that the same memory can be reconstructed under admissible reversal.

\subsection*{2. The audit ladder}
The working audit ladder of this release is
\[
\begin{aligned}
&\text{local triolic constraint} \Longrightarrow \text{HC compensation} \Longrightarrow \text{closure confinement}\\
&\Longrightarrow \text{probe/interface filtration} \Longrightarrow \text{lemma or table compression}\\
&\Longrightarrow \text{readout statement} \Longrightarrow \text{controlled re-expansion}.
\end{aligned}
\]
The final arrow is new in emphasis. It makes explicit that a readout is not the
terminal endpoint of the proof architecture. Instead, it is a temporarily stable
surface from which the architecture should remain recoverable.

\paragraph{Audit Proposition B (endpoints remain reversible).}
A readout that has lost all reversible relation to its closure ladder should be
classified as provisional. The stronger the reversible link to its generating
path, the stronger its status inside the monolith.

\subsection*{3. Probe-audit discipline}
The probe channel receives an additional task in v2.1.0. It is no longer
only a selection or filtration device; it also serves as an audit checkpoint.
Probe-selected modes should leave a downstream trace that remains visible in the
compressed lemma/table/readout layers. This makes the probe sector a bridge
between admissible dynamics and later recoverable proof memory.

\paragraph{Audit Proposition C (probe traces survive compression).}
If a configuration is admissibly selected in the orthogonal probe channel, then
its later compressed representations should preserve a recognizable probe trace.
If no such trace remains, downstream formulations must be marked as reduced,
heuristic, or not yet fully synchronized.

\subsection*{4. Appendix as an audit surface}
This version also sharpens the role of the appendix. Appendix entries, tables,
and named tools now act as \emph{audit surfaces}. Each item should indicate not
only what it abbreviates, but how it can be re-opened into the same structural
class. The appendix therefore becomes more than a memory store or ledger sheet;
it becomes a reversible interface between compressed exposition and detailed
proof ancestry.

\paragraph{Audit Proposition D (appendix entries carry reopening metadata).}
A mature appendix entry should point both backward and forward: backward into
its closure ancestry and forward into an admissible re-expansion channel. This
dual role prevents tables from drifting into detached summaries.

\subsection*{5. Consequence for the next mathematical thickening}
With v2.1.0 the document takes another step toward a larger mathematically
expanded companion. Future lemma growth should now be judged not only by local
correctness, but also by whether it improves auditability across compression
levels. New material is strongest when it reduces the gap between compressed
statement, appendix row, probe trace, and recoverable closure ancestry.

\paragraph{Working rule for subsequent versions.}
Whenever a new mathematical block is added, it should be possible to specify
which closure class licenses it, which compressed interfaces abbreviate it, and
which controlled re-expansion route reconstructs it. Material satisfying all
three demands strengthens the proof architecture in the intended structure-first
sense.



\clearpage
\section*{v2.2.0 Structural Re-Expansion Governance Core}
\addcontentsline{toc}{section}{v2.2.0 Structural Re-Expansion Governance Core}
\begin{quote}
\textbf{Normalization note.} This block is retained as a historical governance precursor. It contributes release and re-expansion discipline, but in the normalized architecture it is not an independent proof center; proof-bearing content is carried by the grouped unfolding phases and their semantic/tensor transport rules.
\end{quote}

This release adds a governance layer to the companion monolith. Earlier stages
made closure transport, probe filtration, auditability, and ledger discipline
explicit. Version~v2.2.0 now asks an additional question: under which
rules may a compressed statement be enlarged again without damaging the closure
class, the appendix interfaces, or the downstream readout status? The answer is
cast as a governance core. Re-expansion is no longer only possible; it is to be
performed under named admissibility rules.

\subsection*{1. Governance as a proof-discipline layer}
A mature structure-first document should not merely accumulate more lemmas. It
should specify which later thickening operations are allowed and which ones are
not. Governance therefore means that every future expansion step is judged by
its compatibility with closure ancestry, probe traces, ledger surfaces, and the
status labels already attached to the claim being enlarged.

\paragraph{Governance Proposition A (admissible thickening preserves status).}
A later elaboration is admissible only if it preserves the structural status of
the compressed claim or improves it in a controlled way. An expansion that
forces a change of admissibility class without explicit declaration should be
regarded as a separate branch rather than a faithful thickening of the original
path.

\paragraph{Proof sketch.}
Compression is useful only because it abbreviates a path already licensed by the
monolith. When a later reader expands the statement again, the same path should
remain visible. If the elaboration silently changes the closure ancestry, probe
selection, or projection label, then it is not a recovery of the earlier path
but a new construction. Governance marks this distinction explicitly.

\subsection*{2. The re-expansion governance ladder}
The practical ladder of this release is
\[
\begin{aligned}
&\text{closure ancestry} \Longrightarrow \text{probe/interface trace} \Longrightarrow \text{ledger or appendix encoding}\\
&\Longrightarrow \text{compressed readout} \Longrightarrow \text{governed re-expansion}\\
&\Longrightarrow \text{status check} \Longrightarrow \text{return to the ledger}.
\end{aligned}
\]
The final return arrow is part of the new emphasis. Re-expansion is not treated
as a one-way release of detail. It must return to the ledger surface in a form
that remains compatible with the proof memory of the monolith.

\paragraph{Governance Proposition B (every governed expansion returns to a ledger surface).}
A mathematically thickened block is fully integrated only when its added detail
can be pushed back into a stable appendix, tool, lemma, or table interface.
Otherwise the expansion remains locally useful but globally insufficiently
synchronized.

\subsection*{3. Probe governance and admissible branch formation}
The probe channel now gains a sharper branching role. Not every elaboration of a
probe-selected mode has to remain in the same local branch, but any branch split
must be declared at the point where the probe trace ceases to determine a unique
downstream continuation. This keeps the search-space logic compatible with the
proof memory logic.

\paragraph{Governance Proposition C (branching must occur at visible trace loss).}
If a later derivation no longer preserves a unique recognizable probe trace,
then a branch declaration is required. Continuing as if the same path remained
in force would obscure which downstream statements are still inherited and which
are newly introduced.

\subsection*{4. Appendix governance and reversible table discipline}
The appendix receives a correspondingly stronger role. Rows, tables, and named
tools are now governed interfaces. They should not only compress results and
carry reopening metadata; they should also state which kinds of later expansion
may legitimately pass through them. In this sense the appendix becomes a modest
control layer for future mathematical thickening.

\paragraph{Governance Proposition D (table rows define admissible recovery windows).}
A mature table row should delimit the window inside which a later elaboration is
still regarded as a recovery of the same path class. Once an expansion leaves
that window, the result should either be relabeled or entered as a new row,
lemma, or branch.

\subsection*{5. Consequence for the next proof cycle}
With v2.2.0 the companion monolith now carries not only closure,
projection, ledger, and audit discipline, but also a clearer rule for how later
mathematical detail may be inserted without dissolving the structure-first
architecture. The next proof cycle should therefore thicken selected paths under
explicit governance rather than by unconstrained local expansion. This is the
cleanest route toward the larger companion target with many more lemmas and
technical proof tools.

\paragraph{Working rule for subsequent versions.}
Whenever a new technical derivation is inserted, it should be possible to say
(i) which closure and probe ancestry it inherits, (ii) which appendix or ledger
surface compresses it, (iii) which re-expansion window legitimizes its added
detail, and (iv) how the resulting thickening returns to a synchronized proof
interface. Material meeting all four demands strengthens the monolith in the
intended disciplined way.


\clearpage
\section*{Mathematical Thickening of the HC/Closure Core (working expansion)}
\addcontentsline{toc}{section}{Mathematical Thickening of the HC/Closure Core (working expansion)}

This block deliberately reverses part of the recent compression trend. It does
not introduce a new meta-layer. Instead it takes one of the central structure
claims --- ground condition $\Rightarrow$ HC $\Rightarrow$ closure --- and writes it
in a denser mathematical form. The purpose is not to declare a final axiomatic
closure theorem, but to make explicit which quantitative objects control the
transition from local triolic structure to confinement and later readout.

\subsection*{1. Working state space and defect variables}
Let $\mathcal H$ be a real or complex Hilbert space with orthogonal splitting
\[
\mathcal H = \mathcal H_A \oplus \mathcal H_B \oplus \mathcal H_C \oplus \mathcal H_R,
\]
where $\mathcal H_A,\mathcal H_B,\mathcal H_C$ carry the three triolic channels and
$\mathcal H_R$ is a residual sector collecting modes not admitted by the intended
closure regime. We write a state as
\[
\Psi = (A,B,C,R), \qquad A\in\mathcal H_A,\ B\in\mathcal H_B,\ C\in\mathcal H_C,
\ R\in\mathcal H_R.
\]
The orthogonality discipline of the triolic block is encoded by the defect
functional
\[
\delta_{\perp}(\Psi)
:= |\langle A,C\rangle|^2 + |\langle B,C\rangle|^2,
\]
and the residual leakage into the non-admissible sector by
\[
\delta_R(\Psi) := \|R\|^2.
\]
The first quantity measures failure of the distinguished $C$-channel to remain
orthogonal to the other two active channels. The second measures failure of the
state to remain confined to the triolic/septinion envelope.

\paragraph{Ground functional.}
A compact way to encode the ground condition is through the non-negative
functional
\[
\Gamma(\Psi) := \delta_{\perp}(\Psi) + \eta\,\delta_R(\Psi), \qquad \eta>0.
\]
The ideal ground regime is characterized by $\Gamma(\Psi)=0$. In particular,
\[
\Gamma(\Psi)=0 \quad \Longrightarrow \quad \langle A,C\rangle=\langle B,C\rangle=0
\quad\text{and}\quad R=0.
\]
Thus the ground condition simultaneously imposes orthogonality and exclusion of
residual leakage.

\subsection*{2. A formal compensator model}
To thicken the HC claim mathematically, introduce the gradient of $\Gamma$ with
respect to the Hilbert structure. Since
\[
\nabla_A \Gamma = 2\langle A,C\rangle C,
\qquad
\nabla_B \Gamma = 2\langle B,C\rangle C,
\]
and
\[
\nabla_C \Gamma = 2\overline{\langle A,C\rangle}A + 2\overline{\langle B,C\rangle}B,
\qquad
\nabla_R \Gamma = 2\eta R,
\]
a natural working expression for the compensator is the feedback field
\[
\mathrm{HC}(\Psi) := -\kappa\,\nabla \Gamma(\Psi), \qquad \kappa>0.
\]
Written componentwise,
\[
\mathrm{HC}(\Psi) =
\Bigl(
-2\kappa\langle A,C\rangle C,
-2\kappa\langle B,C\rangle C,
-2\kappa\overline{\langle A,C\rangle}A -2\kappa\overline{\langle B,C\rangle}B,
-2\kappa\eta R
\Bigr).
\]
This is the simplest explicit model of the structural statement that the
compensator opposes orthogonality defects and suppresses escape into the
residual sector.

\paragraph{Remark.}
The point of this formula is not uniqueness. It is a minimal realization of the
qualitative rule that the HC term points against the defect gradient generated
by the ground functional.

\subsection*{3. Closure as Lyapunov monotonicity}
Assume the effective evolution has the form
\[
\partial_\tau \Psi = F(\Psi) + \mathrm{HC}(\Psi),
\]
where $F$ denotes the non-compensated structural drift. The cleanest closure
statement occurs when the uncontrolled part is tangent or dissipative relative
to the ground functional:
\[
\Re\,\langle \nabla\Gamma(\Psi),F(\Psi)\rangle \le 0.
\tag{T}
\]
Under this hypothesis,
\begin{align*}
\frac{d}{d\tau}\Gamma(\Psi(\tau))
&= \Re\,\langle \nabla\Gamma(\Psi),\partial_\tau\Psi\rangle \\
&= \Re\,\langle \nabla\Gamma(\Psi),F(\Psi)\rangle
   - \kappa\|\nabla\Gamma(\Psi)\|^2 \\
&\le -\kappa\|\nabla\Gamma(\Psi)\|^2 \le 0.
\end{align*}
Hence $\Gamma$ is a Lyapunov functional for the compensated flow.

\paragraph{Lemma 1 (monotone compensated closure).}
Under \textup{(T)}, the map $\tau\mapsto \Gamma(\Psi(\tau))$ is non-increasing. In
particular, the orthogonality defect and residual leakage cannot increase along
the compensated flow.

\paragraph{Corollary 1.}
If $\Gamma(\Psi(0))=0$, then $\Gamma(\Psi(\tau))=0$ for all later $\tau$ for which
the flow exists. Thus the exact ground class is forward invariant.

\paragraph{Corollary 2.}
If $\Gamma(\Psi(0))$ is small, then the trajectory remains trapped in the
sublevel set
\[
\{\Psi\in\mathcal H: \Gamma(\Psi)\le \Gamma(\Psi(0))\}.
\]
This is the first explicit mathematical form of the closure/confinement claim.

\subsection*{4. Septinion confinement as exclusion of the residual sector}
The previous corollaries can be sharpened for the residual component. Since
$\Gamma(\Psi)\ge \eta\|R\|^2$, monotonicity implies
\[
\|R(\tau)\|^2 \le \eta^{-1}\Gamma(\Psi(\tau)) \le \eta^{-1}\Gamma(\Psi(0)).
\]
Therefore a small initial ground defect bounds the entire future residual mass.
In the exact case $\Gamma(\Psi(0))=0$ one obtains $R(\tau)=0$ identically.

\paragraph{Lemma 2 (residual confinement).}
Under \textup{(T)}, the residual sector $\mathcal H_R$ is dynamically suppressed by
the compensator. Exact ground states never leave the triolic/septinion block,
and approximate ground states can leak into $\mathcal H_R$ only at a level
quantitatively controlled by their initial defect.

This is a mathematically sharper version of the earlier verbal statement that
``no component breaks out'' once the HC mechanism is active.

\subsection*{5. First-layer master equation in decomposed form}
The structure-first master equation can now be written in a way that makes the
origin of the principal terms explicit:
\[
\partial_\tau \Psi
= T_{\mathrm{int}}(\Psi)
+ T_{\mathrm{proj}}(\Psi)
- \kappa\nabla\Gamma(\Psi).
\]
Here
\begin{itemize}
    \item $T_{\mathrm{int}}$ collects the intrinsic triolic or septinion-internal
    couplings,
    \item $T_{\mathrm{proj}}$ collects frame-shift or measurement-induced projection
    terms,
    \item $-\kappa\nabla\Gamma$ is the compensating closure term.
\end{itemize}
The point of this decomposition is that the HC contribution is no longer merely
described verbally. It is isolated as the mathematically distinguished term that
penalizes deviation from the admissible structural class.

\paragraph{Lemma 3 (term-wise origin of stabilization).}
Suppose $T_{\mathrm{int}}+T_{\mathrm{proj}}$ satisfies the tangency condition
\textup{(T)}. Then every decrease of the ground functional is generated by the HC
term, while all remaining contributions are either neutral or already
compatible with closure. In that sense the master equation separates
structurally allowed drift from explicitly compensating drift.

\subsection*{6. Probe sector as constrained reduction}
For the probe channel, write two incoming triolic probes as
\[
\Psi_1=(A_1,B_1,C_1,R_1), \qquad \Psi_2=(A_2,B_2,C_2,R_2),
\]
and let $Q$ denote the orthogonal third group used to delimit the search space.
A natural probe defect is then
\[
\Gamma_{\mathrm{probe}}(\Psi_1,\Psi_2;Q)
:= \Gamma(\Psi_1)+\Gamma(\Psi_2)
   + \rho\,\Xi(\Psi_1,\Psi_2;Q), \qquad \rho>0,
\]
where $\Xi$ measures incompatibility with the orthogonal search geometry. For
example, one may take
\[
\Xi(\Psi_1,\Psi_2;Q)
:= |\langle A_1,Q\rangle|^2 + |\langle B_1,Q\rangle|^2
 + |\langle A_2,Q\rangle|^2 + |\langle B_2,Q\rangle|^2,
\]
possibly supplemented by additional coupling terms between the two probes.
Then the same compensated-flow idea yields
\[
\partial_\tau(\Psi_1,\Psi_2)
= F_{\mathrm{probe}}(\Psi_1,\Psi_2;Q)
- \kappa_p\nabla\Gamma_{\mathrm{probe}}(\Psi_1,\Psi_2;Q),
\]
with
\[
\frac{d}{d\tau}\Gamma_{\mathrm{probe}} \le -\kappa_p
\|\nabla\Gamma_{\mathrm{probe}}\|^2
\]
under the corresponding tangency condition.

\paragraph{Lemma 4 (probe-sector reduction principle).}
Under the probe tangency hypothesis, the orthogonal search geometry reduces the
effective probe dynamics to a controlled subspace determined by the sublevel
sets of $\Gamma_{\mathrm{probe}}$. Thus the probe mechanism is mathematically read
as a constrained reduction principle rather than a merely heuristic search
picture.

\subsection*{7. What this thickening achieves and what it does not yet do}
This expansion does three concrete things.
\begin{enumerate}
    \item It replaces the purely verbal HC statement by an explicit compensator
    model $-\kappa\nabla\Gamma$.
    \item It turns closure into a Lyapunov monotonicity statement.
    \item It gives the probe sector a quantitative defect functional that can be
    minimized under orthogonal search constraints.
\end{enumerate}
At the same time, it is still a working expansion rather than a final theorem.
The non-compensated drifts $F$, $T_{\mathrm{int}}$, $T_{\mathrm{proj}}$, and
$F_{\mathrm{probe}}$ have not yet been specialized to a unique physical or number-
theoretic model. That specialization is the next true mathematical step after
this thickening: once the intrinsic drift is fixed more sharply, the above
Lyapunov framework can be tested term-by-term instead of only at the abstract
closure level.


\section*{v2.3.0 Mathematical Deepening Core}
\begin{quote}
\textbf{Normalization note.} This block is retained as a historical deepening precursor. Its mathematical thickening remains part of the document's proof memory, but the normalized architecture reads it through the later unfolding phases rather than as a second autonomous mathematical foundation.
\end{quote}

This version continues the return from governance-style consolidation to
explicit mathematical derivation. The focus is no longer merely that closure can
be encoded by a defect functional, but that the induced closure flow can be
separated into invariant, dissipative, and residual channels. In this way the
HC term becomes a mathematically trackable transport law rather than a purely
qualitative stabilizer.

\subsection*{1. Orthogonal projector formulation of closure}
Let $\mathcal H=\mathcal H_A\oplus\mathcal H_B\oplus\mathcal H_C\oplus\mathcal H_R$
be the structural state space used in the previous thickening step, and let
$\Pi_{\mathrm{cl}}$ denote the orthogonal projector onto the admissible closure
manifold
\[
\mathfrak M_{\mathrm{cl}}:=\{\Psi\in\mathcal H\,:\,\Gamma(\Psi)=0\}.
\]
Whenever $\Gamma$ is locally quadratic near $\mathfrak M_{\mathrm{cl}}$, one may
write, to leading order,
\[
\Gamma(\Psi)\asymp \| (I-\Pi_{\mathrm{cl}})\Psi\|^2,
\]
so that the distance to closure is measured by the defect component
\[
\Psi_{\perp}:=(I-\Pi_{\mathrm{cl}})\Psi,
\qquad
\Psi_{\parallel}:=\Pi_{\mathrm{cl}}\Psi.
\]
The compensated flow
\[
\partial_\tau \Psi = F(\Psi)-\kappa\nabla\Gamma(\Psi)
\]
therefore admits the splitting
\[
\partial_\tau \Psi_{\parallel}=\Pi_{\mathrm{cl}}F(\Psi)+\mathcal N_{\parallel}(\Psi),
\qquad
\partial_\tau \Psi_{\perp}=(I-\Pi_{\mathrm{cl}})F(\Psi)-\kappa\nabla\Gamma(\Psi)+\mathcal N_{\perp}(\Psi),
\]
where $\mathcal N_{\parallel},\mathcal N_{\perp}$ collect higher-order curvature
terms of the closure manifold.

\paragraph{Lemma 1 (normal damping near closure).}
Assume that $\Gamma$ is $C^2$ in a neighborhood of $\mathfrak M_{\mathrm{cl}}$ and that
its Hessian is strictly positive on the normal bundle of
$\mathfrak M_{\mathrm{cl}}$. Then there are constants $c_1,c_2>0$ such that for all
states sufficiently near closure,
\[
c_1\|\Psi_{\perp}\|^2\le \Gamma(\Psi)\le c_2\|\Psi_{\perp}\|^2,
\]
and the HC term damps the normal component at leading order.

\begin{proof}
By the Morse--Bott type expansion of $\Gamma$ around the zero set,
\[
\Gamma(\Psi_{\parallel}+\Psi_{\perp})=
\tfrac12\langle H_{\Gamma}(\Psi_{\parallel})\Psi_{\perp},\Psi_{\perp}\rangle
+O(\|\Psi_{\perp}\|^3),
\]
with $H_{\Gamma}$ positive on normal directions. Hence $\Gamma$ is equivalent to
the squared normal distance, and
$\nabla\Gamma(\Psi)=H_{\Gamma}(\Psi_{\parallel})\Psi_{\perp}+O(\|\Psi_{\perp}\|^2)$.
Therefore the compensator contributes a strictly contracting term in the normal
direction.
\end{proof}

\subsection*{2. Differential inequality for the closure ladder}
Set $E(\tau):=\Gamma(\Psi(\tau))$. Then
\[
\dot E = \langle \nabla\Gamma(\Psi),F(\Psi)\rangle - \kappa\|\nabla\Gamma(\Psi)\|^2.
\]
To make the structural content explicit, split the drift as
\[
F = F_{\mathrm{tan}} + F_{\mathrm{nor}} + F_{\mathrm{res}},
\]
where $F_{\mathrm{tan}}$ is tangent to $\mathfrak M_{\mathrm{cl}}$, $F_{\mathrm{nor}}$ is
controlled by the closure defect, and $F_{\mathrm{res}}$ contains genuinely residual
terms. Then
\[
\dot E
= \langle \nabla\Gamma,F_{\mathrm{nor}}\rangle
 + \langle \nabla\Gamma,F_{\mathrm{res}}\rangle
 - \kappa\|\nabla\Gamma\|^2,
\]
because the tangential part drops out to first order.

\paragraph{Proposition 1 (closure ladder with residual threshold).}
Suppose there exist constants $a<\kappa$ and $b\ge 0$ such that
\[
|\langle \nabla\Gamma,F_{\mathrm{nor}}+F_{\mathrm{res}}\rangle|
\le a\|\nabla\Gamma\|^2 + b\Gamma.
\]
Then for states near closure,
\[
\dot E \le -(\kappa-a)\|\nabla\Gamma\|^2 + bE.
\]
In particular, if $b$ is dominated by the coercivity constant relating
$\|\nabla\Gamma\|^2$ and $E$, then the closure ladder is dissipative and the flow
returns exponentially to the admissible structural sector.

\begin{proof}
Insert the estimate directly into the identity for $\dot E$. The final claim
follows from the local coercivity
$\|\nabla\Gamma\|^2\ge cE$ near $\mathfrak M_{\mathrm{cl}}$, which converts the
inequality into $\dot E\le -\mu E$ for some $\mu>0$.
\end{proof}

This proposition sharpens the earlier Lyapunov statement. Closure is not only
monotone under ideal tangency; it remains stable under a controlled class of
residual perturbations, provided the compensator strength exceeds the residual
threshold.

\subsection*{3. Term-wise master equation refinement}
Write the first-layer master equation in the refined form
\[
\partial_\tau \Psi
= T_{\mathrm{tri}}(\Psi)
+ T_{\perp}(\Psi)
+ T_{\mathrm{meas}}(\Psi)
+ T_{\mathrm{res}}(\Psi)
- \kappa\nabla\Gamma(\Psi).
\]
Here the meaning of the terms is:
\begin{itemize}
    \item $T_{\mathrm{tri}}$: internal triolic transport preserving the intrinsic
    algebraic coupling,
    \item $T_{\perp}$: orthogonality-restoring channel induced by the requirement
    that one branch remains perpendicular to the other two,
    \item $T_{\mathrm{meas}}$: frame-shift and measurement transport,
    \item $T_{\mathrm{res}}$: unresolved residual transport not yet assigned to a
    unique physical model.
\end{itemize}
The closure condition can then be tested term by term through the scalar
pairings
\[
\sigma_{\mathrm{tri}}:=\langle \nabla\Gamma,T_{\mathrm{tri}}\rangle,
\quad
\sigma_{\perp}:=\langle \nabla\Gamma,T_{\perp}\rangle,
\quad
\sigma_{\mathrm{meas}}:=\langle \nabla\Gamma,T_{\mathrm{meas}}\rangle,
\quad
\sigma_{\mathrm{res}}:=\langle \nabla\Gamma,T_{\mathrm{res}}\rangle.
\]
Thus
\[
\dot E = \sigma_{\mathrm{tri}}+\sigma_{\perp}+\sigma_{\mathrm{meas}}+\sigma_{\mathrm{res}}
-\kappa\|\nabla\Gamma\|^2.
\]
The mathematical advantage is that every structural claim can now be tied to one
of these four channels. For instance, a perfect intrinsic closure sector is
characterized by $\sigma_{\mathrm{tri}}=0$, while measurement-induced distortion is
captured by the sign and size of $\sigma_{\mathrm{meas}}$.

\paragraph{Corollary 2 (term-wise admissibility test).}
If each non-compensated channel satisfies an estimate of the form
$|\sigma_j|\le a_j\|\nabla\Gamma\|^2+b_jE$ with $\sum_j a_j<\kappa$, then the
master flow is closure-admissible. In this sense the first-layer master
equation acquires a concrete mathematical admissibility criterion instead of a
purely narrative one.

\subsection*{4. Probe reduction with quantitative search radius}
For the probe functional
\[
\Gamma_{\mathrm{probe}}(\Psi_1,\Psi_2;Q)
=\Gamma(\Psi_1)+\Gamma(\Psi_2)+\rho\,\Xi(\Psi_1,\Psi_2;Q),
\]
define the admissible search radius
\[
\mathcal R_\varepsilon(Q):=
\{(\Psi_1,\Psi_2):\Gamma_{\mathrm{probe}}(\Psi_1,\Psi_2;Q)\le \varepsilon\}.
\]
This produces a genuinely quantitative reading of the earlier search-space
language: the probe channel is restricted to a sublevel tube around the
orthogonally compatible sector.

\paragraph{Proposition 3 (quantitative probe reduction).}
Assume the probe drift satisfies the analogue of the closure-ladder estimate.
Then for every sufficiently small $\varepsilon>0$ the set
$\mathcal R_\varepsilon(Q)$ is forward invariant under the compensated probe flow,
provided the initial probe state lies in $\mathcal R_\varepsilon(Q)$. Hence the
probe search does not merely heuristically prefer a smaller space; it is
mathematically confined to a controlled admissible tube.

\begin{proof}
Apply the same differential-inequality argument to
\[
E_{\mathrm{probe}}(\tau):=\Gamma_{\mathrm{probe}}(\Psi_1(\tau),\Psi_2(\tau);Q).
\]
If $\dot E_{\mathrm{probe}}\le -\mu E_{\mathrm{probe}}$ or, more generally, $\dot E_{\mathrm{probe}}\le 0$ on the boundary level set $E_{\mathrm{probe}}=\varepsilon$, trajectories cannot leave the sublevel region.
\end{proof}

\subsection*{5. Consequence for the next development step}
The present deepening changes the status of the HC/closure block in two ways.
First, closure is now expressed through projector splitting and coercive defect
control. Second, the master and probe channels both admit explicit admissibility
tests in terms of scalar defect pairings. The next natural step is therefore no
longer another meta-level consolidation, but the specialization of
$T_{\mathrm{tri}},T_{\perp},T_{\mathrm{meas}},T_{\mathrm{res}}$ and of the probe drift to
more definite structural or physical models. That is the point at which the
current structure-first proof begins to turn into a truly term-resolved
mathematical derivation.


\section*{v2.4.0 Term-Resolved Channel Specialization Core}
\addcontentsline{toc}{section}{v2.4.0 Term-Resolved Channel Specialization Core}
\begin{quote}
\textbf{Normalization note.} This block is retained as a historical specialization precursor. In the normalized architecture its active content belongs primarily to the projection/translation and field/probe realization layers; it should therefore be read as a later historical expansion, not as a competing core.
\end{quote}

The explicit task of v\docversion\ is to pass from the admissibility language of
v2.3.0 to a first term-resolved structural model. The point is not yet to claim
final physics. Rather, the point is to assign mathematically disciplined roles to
\(T_{\mathrm{tri}},T_{\perp},T_{\mathrm{meas}},T_{\mathrm{res}}\) so that later field equations,
appendix tables, and probe readouts can all be tested against the same channel map.
The guiding principle is that every term must be readable simultaneously in four
ways: as a structural transport rule, as a closure-defect source or sink, as a
probe-interface contribution, and as a future physical interpretation slot.

\subsection*{1. Channel factorization principle}
Let the first-layer master equation be written as
\[
\partial_\tau \Psi
= T_{\mathrm{tri}}(\Psi)
+ T_{\perp}(\Psi)
+ T_{\mathrm{meas}}(\Psi)
+ T_{\mathrm{res}}(\Psi)
- \kappa \nabla\Gamma(\Psi).
\]
In v\docversion\ we impose the factorization rule
\[
T_j = A_j + B_j + C_j,
\qquad j\in\{\mathrm{tri},\perp,\mathrm{meas},\mathrm{res}\},
\]
where:
\begin{itemize}
    \item $A_j$ is the closure-tangent component,
    \item $B_j$ is the explicitly defect-producing normal component,
    \item $C_j$ is a bookkeeping remainder reserved for later model refinement.
\end{itemize}
This factorization is not claimed to be unique. Its function is instead to force
each channel into a common proof grammar: tangent transport, normal defect, and
controlled unresolved remainder.

\paragraph{Definition 1 (channel-adapted splitting).}
Given the projector pair $(P_{\mathrm{tan}},P_{\mathrm{nor}})$ induced by the closure
manifold, define
\[
A_j:=P_{\mathrm{tan}}T_j,
\qquad
B_j:=P_{\mathrm{nor}}T_j,
\qquad
C_j:=0
\]
at the first structural level. More general splittings with nonzero $C_j$ are
allowed when one wants to isolate scale-separated or frame-shifted contributions.
At the present stage, however, the projector-induced splitting is the canonical
choice because it minimizes extra assumptions.

\subsection*{2. Canonical specialization of the four channels}
The structural role of the four channels is now fixed as follows.

\paragraph{Triolic channel.}
The term $T_{\mathrm{tri}}$ is the intrinsic coupling transport generated by the
triolen architecture itself. Its preferred structural property is near tangency:
\[
P_{\mathrm{nor}}T_{\mathrm{tri}} = O(\Gamma).
\]
This expresses that the intrinsic triadic coupling should preserve admissibility to
leading order and create only higher-order closure defects.

\paragraph{Orthogonality channel.}
The term $T_{\perp}$ enforces the standing project constraint that one branch is
perpendicular to the other two. Its preferred structural property is restorative:
\[
\langle \nabla\Gamma, T_{\perp}\rangle \le -\lambda_{\perp}\Gamma
\quad\text{for some }\lambda_{\perp}>0
\]
in the near-closure regime. Thus $T_{\perp}$ is not merely admissible; it is a
coercive correction channel.

\paragraph{Measurement channel.}
The term $T_{\mathrm{meas}}$ encodes the phase-space shift between the measurement
frame and the underlying structural dynamics. It is therefore not expected to be
sign-definite. Instead it is controlled by a relative bound of the form
\[
|\langle \nabla\Gamma, T_{\mathrm{meas}}\rangle|
\le a_{\mathrm{meas}}\|\nabla\Gamma\|^2 + b_{\mathrm{meas}}\Gamma.
\]
This formalizes the standing project rule that measurements do not access the raw
structure directly, but only through a shifted interface.

\paragraph{Residual channel.}
The term $T_{\mathrm{res}}$ is the disciplined storage slot for all contributions not
yet promoted to a stable structural or physical interpretation. In v\docversion\ it is no longer allowed to function as a dumping ground. Instead one requires
\[
T_{\mathrm{res}} = \sum_{m=1}^M T_{\mathrm{res}}^{(m)}
\]
with every summand tagged by its origin class, such as approximation residue,
frame-conversion remainder, truncation error, or phenomenological placeholder.
This turns residuality into an auditable notion.

\subsection*{3. Channel ledger and admissibility matrix}
To make the specialization operational, v2.4.0 introduces the channel
ledger matrix
\[
\Lambda_{\mathrm{ch}}(\Psi)=
\begin{pmatrix}
\ell_{\mathrm{tri}} & \delta_{\mathrm{tri}} & r_{\mathrm{tri}}\\
\ell_{\perp} & \delta_{\perp} & r_{\perp}\\
\ell_{\mathrm{meas}} & \delta_{\mathrm{meas}} & r_{\mathrm{meas}}\\
\ell_{\mathrm{res}} & \delta_{\mathrm{res}} & r_{\mathrm{res}}
\end{pmatrix},
\]
where
\[
\ell_j := \|P_{\mathrm{tan}}T_j\|,
\qquad
